矩阵微积分最常见的困难，是不知道结果``应该长成什么形状''，具体计算反而居于其后。不同资料还可能把同一 Jacobian 写成互相转置的形式，使公式表看起来彼此冲突。

解决办法是先回到导数的第一性定义：导数是输入微小变化到输出一阶变化之间的线性映射。

\begin{enumerate}
\tightlist
\item
  导数是最佳线性映射：从 Fréchet 导数出发，统一标量、向量与矩阵输入。
\item
  形状先于公式，布局必须声明：选择``输出按行、输入按列''的 Jacobian 约定，并解释梯度、JVP、VJP 与其他资料的转置关系。
\end{enumerate}

后续推导会尽量先写微分或线性作用，再把它读成梯度。这样，即使换了一套布局约定，数学映射本身也不会改变。

\subsection{两个约定对比一目了然}

设 \(y=Ax\)，\(A\in\R^{2\times3}\)，\(x\in\R^3\)，\(y\in\R^2\)。在``输出行、输入列''约定下，Jacobian 是 \(A\) 本身，形状 \(2\times3\)。在转置约定下，Jacobian 是 \(A^{\trans}\)，形状 \(3\times2\)。两种表示都正确表达了``\(dy=A\,dx\)''这个线性映射——区别只在于矩阵乘法中因子写在左边还是右边。在作者统一约定下，所有后续公式自洽；在两种约定间直接搬运公式而不调整乘法顺序，才会出现矛盾。

\subsection{怎么读这一章}

先暂时放下矩阵求导公式表。第一篇把导数还原成``输入扰动怎样线性地传到输出''，第二篇再讨论这张线性映射写成矩阵时为何会出现布局差异。每得到一个结果，先核对输入、输出和扰动的形状，再核对数值。

\subsection{本章篇目}

以下篇目按概念依赖排列。

\begin{itemize}
\tightlist
\item \hyperref[sec:10-Matrix-Calculus/01-Derivative-Conventions/01]{``导数是最佳线性映射''}；
\item \hyperref[sec:10-Matrix-Calculus/01-Derivative-Conventions/02]{``形状先于公式，布局必须声明''}。
\end{itemize}
